Work first conditionally on \(A\) and \((f_i)_{i\in\mathcal V_n}\).  For every \(i\in\mathcal V_n\), define the proof-intermediate function
\[
g_i(x):=\frac{f_i(1,x)+f_i(0,x)}{2},\qquad x\in[0,1].
\]
Fix \(i\) with \(N_i=m>0\).  By \(\texttt{ass:bernoulli-design}\), conditional on \(A\) and all preassignment variables, \(M_i\sim {\rm Binomial}(m,1/2)\), and \(Z_i\) is independent of \(M_i\).  The conditional assignment law does not depend on the latent types, so the same assertions hold after conditioning only on \(A\) and \((f_j)_{j\in\mathcal V_n}\), by iterated conditional expectation.  Hence \(\texttt{ass:anonymous-interference}\) and \(\texttt{def:li-wager-target}\) give the unit-level contribution
\[
\begin{aligned}
\beta_i
&:=4\,\mathbb E_Z\!\left[Y_i(Z)(M_i-m/2)\mid A,(f_j)_{j\in\mathcal V_n}\right]\\
&=4m\,\mathbb E_Z\!\left[g_i(1/2+V_i)V_i\mid A,(f_j)_{j\in\mathcal V_n}\right].
\end{aligned}
\]
If \(m=0\), then \(M_i-N_i/2=0\), so the same definition gives \(\beta_i=0\).  Consequently,
\[
\bar\tau_{\mathrm{IND},n}=\frac1n\sum_{i\in\mathcal V_n}\beta_i.
\]

For \(m>0\), let \(X_1,\ldots,X_m\) be independent variables taking the values \(-1/2\) and \(1/2\) with equal probability, and put \(S_m=\sum_{j=1}^mX_j\).  Conditionally on the quantities fixed above, \(V_i\) has the law of \(S_m/m\).  Independence and centering, together with \(\mathbb E X_j^2=1/4\) and \(\mathbb E X_j^4=1/16\), yield
\[
\mathbb E S_m^2=\frac m4,
\qquad
\mathbb E S_m^3=0,
\]
and, because only the one-index and two-pair terms survive in the fourth-power expansion,
\[
\mathbb E S_m^4
=\sum_{j=1}^m\mathbb E X_j^4
+6\sum_{1\le j<k\le m}\mathbb E X_j^2\mathbb E X_k^2
=\frac{3m^2-2m}{16}.
\]
Therefore
\[
\mathbb E V_i=0,
\qquad
\mathbb E V_i^2=\frac1{4m},
\qquad
\mathbb E V_i^3=0,
\qquad
\mathbb E V_i^4=\frac{3m^2-2m}{16m^4}\le\frac3{16m^2}.
\]

By \(\texttt{ass:third-derivative-bound}\), the averaging defining \(g_i\) satisfies
\[
\sup_{x\in[0,1]}|g_i^{(3)}(x)|\le L.
\]
Repeated integration of \(g_i^{(3)}\) along the segment from \(1/2\) to \(1/2+v\), which lies in \([0,1]\) for every realized \(v=V_i\), gives
\[
g_i(1/2+v)
=g_i(1/2)+g_i'(1/2)v+\frac12g_i''(1/2)v^2+R_i(v),
\qquad
|R_i(v)|\le\frac L6|v|^3.
\]
Multiplying by \(v\), taking conditional expectations, and using the four moment identities shows that the constant and quadratic terms vanish, while \(4m\mathbb E V_i^2=1\).  Hence
\[
\beta_i-g_i'(1/2)
=4m\,\mathbb E\!\left[V_iR_i(V_i)\mid A,(f_j)_{j\in\mathcal V_n}\right],
\]
and thus
\[
|\beta_i-g_i'(1/2)|
\le \frac{4mL}{6}\mathbb E V_i^4
\le \frac{4mL}{6}\frac3{16m^2}
=\frac{L}{8m}.
\]
By \(\texttt{def:target}\), the contribution of unit \(i\) to \(\theta_n\) is \(g_i'(1/2)\).  If \(N_i=0\), then \(\texttt{ass:response-envelope}\) implies
\[
|g_i'(1/2)|
\le\frac{|\partial_x f_i(1,1/2)|+|\partial_x f_i(0,1/2)|}{2}
\le B.
\]
Separating nonisolated and isolated vertices and applying the triangle inequality now proves
\[
\begin{aligned}
|\bar\tau_{\mathrm{IND},n}-\theta_n|
&\le \frac1n\sum_{i:N_i>0}|\beta_i-g_i'(1/2)|
+\frac1n\sum_{i:N_i=0}|g_i'(1/2)|\\
&\le \frac{L}{8n}\sum_{i:N_i>0}\frac1{N_i}
+\frac Bn\sum_{i\in\mathcal V_n}\mathbf1\{N_i=0\}.
\end{aligned}
\]

It remains to establish the asserted asymptotic comparison.  Consider a triangular-array sequence in the Li--Wager Theorem~6 regime, hence satisfying all hypotheses of that theorem.  The uniform latent-type marginal is the one recorded in the repaired Li--Wager sampling condition \(\texttt{ass:unit-sampling}\), and \(\texttt{lem:li-wager-theorem-six-regime}\) gives
\[
\liminf_{n\to\infty}\frac{\log\rho_n}{\log n}>-\frac12.
\]
Choose \(\varepsilon>0\) strictly smaller than the gap between this lower limit and \(-1/2\).  Then, for all sufficiently large \(n\),
\[
\rho_n\ge n^{-1/2+\varepsilon}.
\]
By the independence and uniform marginal in \(\texttt{ass:unit-sampling}\), and the conditional edge independence in \(\texttt{ass:graphon-draw}\), the pairs governing edges incident to \(i\) are independent after conditioning on \(U_i=u\).  Thus
\[
N_i\mid U_i=u\sim {\rm Binomial}(n-1,p_n(u)),
\qquad
p_n(u)=\rho_n\int_0^1W(u,v)\,dv.
\]
The displayed parameter belongs to \([0,1]\) by \(\texttt{ass:edge-probability}\) and \(\texttt{ass:kernel-bound}\).  By \(\texttt{ass:degree-regularity}\), for Lebesgue-almost every \(u\),
\[
\mu_i(u):=\mathbb E[N_i\mid U_i=u]
=(n-1)p_n(u)
\ge\frac{(n-1)\rho_n}{\kappa}.
\]
Because every \(U_i\) is uniform, all realized latent types lie in this full-measure set almost surely.

We derive the required lower-tail estimate directly.  If \(S\sim{\rm Binomial}(m,p)\), \(\mu=mp\), and \(t=\log2\), exponential Markov's inequality and \(1-y\le e^{-y}\) give
\[
\begin{aligned}
\mathbb P(S\le\mu/2)
&\le e^{t\mu/2}\mathbb E e^{-tS}
=e^{t\mu/2}(1-p+pe^{-t})^m\\
&\le\exp\!\left\{\frac{t\mu}{2}-\mu(1-e^{-t})\right\}
=\exp\!\left\{-\frac{1-\log2}{2}\mu\right\}
\le e^{-\mu/8},
\end{aligned}
\]
where the last inequality uses \((1-\log2)/2\ge1/8\).  Apply this estimate conditionally on each \(U_i\), use the preceding lower bound for \(\mu_i(U_i)\), and take a union bound.  For
\[
\mathcal E_n
:=\left\{\min_{i\in\mathcal V_n}N_i\ge\frac{(n-1)\rho_n}{2\kappa}\right\},
\]
we obtain
\[
\mathbb P(\mathcal E_n^c)
\le n\exp\!\left\{-\frac{(n-1)\rho_n}{8\kappa}\right\}
\longrightarrow0,
\]
because \(n\rho_n\ge n^{1/2+\varepsilon}\) eventually.  On \(\mathcal E_n\) there are no isolates, and the finite-sample inequality already proved yields
\[
\frac{|\bar\tau_{\mathrm{IND},n}-\theta_n|}{\sqrt{\rho_n}}
\le\frac{L\kappa}{4(n-1)\rho_n^{3/2}}.
\]
Moreover,
\[
(n-1)\rho_n^{3/2}
\ge(1-o(1))n^{1/4+3\varepsilon/2}
\longrightarrow\infty.
\]
The radii \(L\) and \(\kappa\) are fixed class constants.  Combining the last three displays proves
\[
\bar\tau_{\mathrm{IND},n}-\theta_n=o_P(\sqrt{\rho_n}).
\]
Finally, \(\texttt{lem:li-wager-theorem-six-regime}\) records that the published PC-balancing Gaussian limit is centered at \(\bar\tau_{\mathrm{IND},n}\) on the same \(\sqrt{\rho_n}\) scale.  The last display therefore makes that result an asymptotic benchmark for \(\theta_n\) along every nondegenerate sequence satisfying all of Theorem~6's hypotheses.  The finite-sample functionals remain distinct by \(\texttt{def:target}\) and \(\texttt{def:li-wager-target}\), exactly as the degree-exact inequality records.